\clearpage
\section{Diagrama de transição}
\label{sec:diagramas}

Esta seção apresenta a representação gráfica dos Autômatos
Finitos Determinísticos definidos na Seção~\ref{sec:modelagem}. Os diagramas
demonstram os estados, as transições e os estados finais de
cada reconhecedor.

Em todos os diagramas, as setas indicam o sentido da leitura
dos caracteres. O estado inicial é indicado por uma seta
externa, enquanto os estados finais são representados por
círculos duplos.

As transições não representadas explicitamente direcionam
a leitura para o estado de rejeição $q_R$, que possui
transição para si mesmo para qualquer símbolo do alfabeto.



\tikzset{
    estado/.style={
        state,
        minimum size=7.5mm,
        inner sep=1pt,
        font=\footnotesize
    },
    inicial/.style={
        estado,
        initial,
        initial text={}
    },
    final/.style={
        estado,
        accepting
    },
    rejeicao/.style={
        estado,
        fill=red!10
    },
    transicao/.style={
        ->,
        >=Stealth,
        semithick,
        font=\footnotesize
    }
}


\subsection{AFD 01 -- Padrão Nacional Único (PNU)}

O primeiro autômato reconhece placas compostas por três
letras seguidas de quatro algarismos, permitindo a presença
opcional de um hífen após a terceira letra.

O diagrama apresenta dois caminhos possíveis a partir
do estado $q_3$: um caminho direto para a leitura dos
algarismos, quando não existe hífen, e outro que passa
pelo estado $q_4$, quando o hífen está presente.

\begin{figure}[ht]
\centering
\resizebox{\linewidth}{!}{%
\begin{tikzpicture}[node distance=0.75cm]

    \node[inicial] (q0) {$q_0$};
    \node[estado, right=of q0] (q1) {$q_1$};
    \node[estado, right=of q1] (q2) {$q_2$};
    \node[estado, right=of q2] (q3) {$q_3$};

    \node[estado, below=0.95cm of q3] (q4) {$q_4$};

    \node[estado, right=1.1cm of q3] (q5) {$q_5$};
    \node[estado, right=of q5] (q6) {$q_6$};
    \node[estado, right=of q6] (q7) {$q_7$};
    \node[final, right=of q7] (q8) {$q_8$};

    \node[rejeicao, below=1.55cm of q5] (qr) {$q_R$};

    \path[transicao]
        (q0) edge node[above] {$\Lt$} (q1)
        (q1) edge node[above] {$\Lt$} (q2)
        (q2) edge node[above] {$\Lt$} (q3)

        (q3) edge node[above] {$\Nm$} (q5)
        (q3) edge node[left] {$\Hf$} (q4)
        (q4) edge node[below right] {$\Nm$} (q5)

        (q5) edge node[above] {$\Nm$} (q6)
        (q6) edge node[above] {$\Nm$} (q7)
        (q7) edge node[above] {$\Nm$} (q8)

        (qr) edge[loop below]
        node[below] {$\Sigma$} (qr);

\end{tikzpicture}%
}
\caption{Diagrama de transição do AFD 01 -- PNU (transições omitidas direcionam para $q_R$).}
\label{fig:afd-pnu}
\end{figure}

O estado final é $q_8$, alcançado após a leitura das
três letras e dos quatro algarismos. A presença do hífen
é opcional, sendo permitida somente após a terceira letra.


\subsection{AFD 02 -- Padrão Mercosul}

O segundo autômato reconhece placas no formato LLLNLNN,
composto por três letras, um algarismo, uma letra e dois
algarismos.

Diferentemente do primeiro autômato, este reconhecedor
não permite a utilização de hífen em nenhuma posição.

\begin{figure}[ht]
\centering
\resizebox{\linewidth}{!}{%
\begin{tikzpicture}[node distance=0.8cm]

    \node[inicial] (q0) {$q_0$};
    \node[estado, right=of q0] (q1) {$q_1$};
    \node[estado, right=of q1] (q2) {$q_2$};
    \node[estado, right=of q2] (q3) {$q_3$};
    \node[estado, right=of q3] (q4) {$q_4$};
    \node[estado, right=of q4] (q5) {$q_5$};
    \node[estado, right=of q5] (q6) {$q_6$};
    \node[final, right=of q6] (q7) {$q_7$};

    \node[rejeicao, below=1.4cm of q4] (qr) {$q_R$};

    \path[transicao]
        (q0) edge node[above] {$\Lt$} (q1)
        (q1) edge node[above] {$\Lt$} (q2)
        (q2) edge node[above] {$\Lt$} (q3)
        (q3) edge node[above] {$\Nm$} (q4)
        (q4) edge node[above] {$\Lt$} (q5)
        (q5) edge node[above] {$\Nm$} (q6)
        (q6) edge node[above] {$\Nm$} (q7)

        (qr) edge[loop below]
        node[below] {$\Sigma$} (qr);

\end{tikzpicture}%
}
\caption{Diagrama de transição do AFD 02 -- Mercosul (transições omitidas direcionam para $q_R$).}
\label{fig:afd-mercosul}
\end{figure}

O estado final $q_7$ indica que toda a sequência exigida
foi reconhecida corretamente. Qualquer caractere fora
da sequência esperada conduz ao estado de rejeição $q_R$.

\subsection{AFD 03 -- Mercosul convertido}

O terceiro autômato é uma extensão do reconhecedor
Mercosul. Sua estrutura permanece LLLNLNN, porém a
quinta posição da placa aceita exclusivamente letras
pertencentes ao conjunto $\LAJ$ (A até J).

Por esse motivo, a única diferença em relação ao AFD 02
está na transição que parte de $q_4$.

\begin{figure}[ht]
\centering
\resizebox{\linewidth}{!}{%
\begin{tikzpicture}[node distance=0.8cm]

    \node[inicial] (q0) {$q_0$};
    \node[estado, right=of q0] (q1) {$q_1$};
    \node[estado, right=of q1] (q2) {$q_2$};
    \node[estado, right=of q2] (q3) {$q_3$};
    \node[estado, right=of q3] (q4) {$q_4$};
    \node[estado, right=of q4] (q5) {$q_5$};
    \node[estado, right=of q5] (q6) {$q_6$};
    \node[final, right=of q6] (q7) {$q_7$};

    \node[rejeicao, below=1.4cm of q4] (qr) {$q_R$};

    \path[transicao]
        (q0) edge node[above] {$\Lt$} (q1)
        (q1) edge node[above] {$\Lt$} (q2)
        (q2) edge node[above] {$\Lt$} (q3)
        (q3) edge node[above] {$\Nm$} (q4)
        (q4) edge node[above] {$\LAJ$} (q5)
        (q5) edge node[above] {$\Nm$} (q6)
        (q6) edge node[above] {$\Nm$} (q7)

        (qr) edge[loop below]
        node[below] {$\Sigma$} (qr);

\end{tikzpicture}%
}
\caption{Diagrama de transição do AFD 03 -- Mercosul convertido (transições omitidas direcionam para $q_R$).}
\label{fig:afd-convertido}
\end{figure}

Nesse reconhecedor, caso seja lida uma letra fora do
conjunto $\LAJ$ na quinta posição, a entrada será
direcionada ao estado de rejeição $q_R$.

\subsection{Considerações sobre os diagramas}

Os três diagramas representam visualmente as funções
de transição especificadas na modelagem matemática.

O AFD 01 diferencia-se por permitir duas possibilidades
de leitura após a terceira letra, enquanto os AFDs 02
e 03 possuem uma sequência linear de transições.

O AFD 03 mantém a estrutura do AFD 02, restringindo
apenas o conjunto de caracteres aceitos na quinta posição.

Dessa forma, os autômatos permitem reconhecer os padrões
de placas definidos no projeto, garantindo que somente
sequências pertencentes às linguagens especificadas
sejam aceitas.